\section{Exceptional nonopposition from an exterior-square pairing}

\subsection{BR-123-WEDGE: a root-pair obstruction}
We continue the original question in cote 123, physical page 22,
footer 21: whether all choices of containing Borels for tori in
Kostant apposition are opposite.  We give complete counterexamples
in types $F_4$ and $E_8$.  The argument uses the adjoint representation
alone, and does not assume a table of exceptional representation
dimensions or a numerical diagonalization.

The known structural inputs are Kostant (2010), Theorems 1.2--1.3,
p.~6, and the principal-element paragraph on p.~7; and Milne (2017),
Proposition 22.17 and Theorems 22.18--22.19, pp.~467--468,
and Proposition 22.41, p.~477.  The latter give the highest-weight
line and complete reducibility in characteristic zero.  The explicit
Euclidean root sets are in Milne, \S24(l), pp.~539--542.  The finite
root calculations actually needed below are displayed, rather than
replaced by a claim that these sources contain the counterexamples.

Let $\mathfrak g=\operatorname{Lie}G$, and let $K$ denote its Killing
form in this section only.  It induces the invariant bilinear form
\[
 K_2(a\wedge b,u\wedge v)
       =K(a,u)K(b,v)-K(a,v)K(b,u)
       \quad\hbox{on }\bigwedge^2\mathfrak g.
 \tag{EX.1}
\]
For a maximal torus with chosen simple roots, suppose that distinct
roots $\theta,\delta$ satisfy the following explicit conditions:
$\theta+\alpha_i$ is never a root; $\delta+\alpha_i$ is a root
only for $i=j$, when it equals $\theta$; and $-1$ belongs to the
Weyl group.  Choose nonzero root vectors and put
\[
                         z=e_\theta\wedge e_\delta.
\]
Every simple positive root operator kills $z$: it kills the first
factor, and its only possible nonzero action on the second produces
a multiple of the first factor.  The simple root operators generate
the positive nilpotent Lie algebra, so the whole unipotent radical
fixes $z$.  It is also a torus eigenvector.  It therefore spans the
highest line of its cyclic subrepresentation $V=\langle Gz\rangle$.
By complete reducibility and Proposition 22.17 this subrepresentation
is irreducible: in a direct-sum decomposition, every summand receiving
the cyclic vector has that same highest weight, and the cyclic
diagonal copy has one-dimensional highest-weight space.

A representative of $-1$ sends $z$ to a nonzero scalar multiple of
$e_{-\theta}\wedge e_{-\delta}$.  This vector belongs to $V$ and
has nonzero $K_2$-pairing with $z$, because opposite root spaces pair
nondegenerately and the two cross-pairings vanish.  Hence $K_2|_V$
is nonzero; its radical is invariant, so irreducibility makes the
restriction nondegenerate.  It follows also directly that highest
lines for opposite Borels in $V$ pair nontrivially: conjugate the
opposite Borels simultaneously to the chosen positive and negative
Borels, whose highest lines are the two just displayed, and use
invariance and uniqueness of the highest line.

Now let $(T_0,T_1)$ be an apposition pair, with principal element
$c\in T_0\cap N_G(T_1)$ of order $h$, acting on the second root
system as a Coxeter element $C$.  Choose the positive root system
of $T_0$ so that
\[
 \alpha_i(c)=\zeta=e^{2\pi i/h},\qquad
 \alpha(c)=\zeta^{\operatorname{ht}(\alpha)}.
 \tag{EX.2}
\]
Suppose that the unordered pairs $(\alpha,\beta)$ in the first
root system and $(\mu,C^d\mu)$ in the second are Weyl translates
of $(\theta,\delta)$, and that
\[
             h\mid d\bigl(\operatorname{ht}(\alpha)
                            -\operatorname{ht}(\beta)\bigr).
 \tag{EX.3}
\]
The vector $z_0=e_\alpha\wedge e_\beta$ is then a highest-line
vector for some Borel $B_0$ containing $T_0$.  If
$u\in\mathfrak g_\mu(T_1)$ is nonzero, then
$z_1=u\wedge\operatorname{Ad}(c^d)u$ is nonzero and is a
highest-line vector for some Borel $B_1$ containing $T_1$.
Both lie in the same subrepresentation $V$: transport the highest
wedge to the other torus by conjugation and then use the stated
Weyl translates.

Write $a=K(e_\alpha,u)$ and $b=K(e_\beta,u)$.  Invariance gives
\[
 K(e_\alpha,\operatorname{Ad}(c^d)u)
       =\zeta^{-d\operatorname{ht}(\alpha)}a,
\]
and similarly for $\beta$.  Consequently
\[
 K_2(z_0,z_1)=ab\bigl(
       \zeta^{-d\operatorname{ht}(\beta)}
       -\zeta^{-d\operatorname{ht}(\alpha)}\bigr)=0.
 \tag{EX.4}
\]
This conclusion does not require either $a$ or $b$ to be nonzero.
The highest-line test proved above shows that $B_0,B_1$ cannot be
opposite.  This proves the obstruction.

For its applications, one can first choose the explicitly ordered
Coxeter element on $T_1$ and a lift $c$.  Kostant's p.~7 statement
makes it principal and regular; its centralizer identity component
is $T_0$.  Conjugating this principal element gives (EX.2), while
transporting the chosen root coordinates on $T_1$.  Thus prescribing
the Coxeter word on the second root system and the principal
coordinates on the first is legitimate.  CC-APG01 then transports
one counterexample to every ordered apposition pair.

\subsection{BR-123-OPPF4: a counterexample in type \texorpdfstring{$F_4$}{F4}}
Use the Euclidean root system
\[
 \{\pm e_i,\ \pm e_i\pm e_j,\
             \tfrac12(\pm e_1\pm e_2\pm e_3\pm e_4)\}
\]
with the simple roots
\[
 a_1=e_2-e_3,\quad a_2=e_3-e_4,\quad a_3=e_4,\quad
 a_4=\tfrac12(e_1-e_2-e_3-e_4).
 \tag{EX.5}
\]
The root-coroot Cartan matrix, with entry
$A_{ij}=\langle a_i,a_j^\vee\rangle$, is
\[
 A=\begin{pmatrix}
 2&-1&0&0\\-1&2&-2&0\\0&-1&2&-1\\0&0&-1&2
 \end{pmatrix}.
\]
Put $\theta=e_1+e_2=2a_1+3a_2+4a_3+2a_4$ and
$\delta=e_1+e_3=\theta-a_1$.  From the displayed root set,
$\theta+a_i$ is never a root, and $\delta+a_i$ is a root only
for $i=1$, giving $\theta$.  The element $-1$ is a Weyl element:
multiply the four reflections in the mutually orthogonal short
roots $e_i$.  Thus the highest-wedge hypotheses hold.

On the first torus take
\[
                 \alpha=e_2+e_3,
                 \qquad\beta=\theta.
\]
Starting with $(\theta,\delta)$, apply the simple reflections
$s_2,s_3,s_4$ in that order.  They fix $\theta$ and send
\[
 e_1+e_3\longmapsto e_1+e_4
       \longmapsto e_1-e_4\longmapsto e_2+e_3.
\]
This is the required unordered-pair Weyl certificate.  Its heights
are $\operatorname{ht}(\alpha)=5$ and
$\operatorname{ht}(\beta)=11$; equivalently the height functional
on Euclidean coordinates is $(8,3,2,1)$.

On the second root system choose $C=s_1s_2s_3s_4$, with the rightmost
reflection applied first, and take $\mu=\theta$.  Direct reflection
gives
\[
 C(e_1+e_2)=e_1+e_3,\qquad C^2(e_1+e_2)=e_1+e_4.
\]
The pair $(\mu,C^2\mu)$ is the translate of
$(\theta,\delta)$ by $s_2$.  Since $h=12$ and
$2(11-5)=12$, (EX.3) holds with $d=2$.  Formula (EX.4) supplies
nonopposite containing Borels.  This proves the claimed failure
of the all-choice assertion in type $F_4$.  $\square$

\subsection{BR-123-OPPE8: a counterexample in type \texorpdfstring{$E_8$}{E8}}
Use the roots $\pm e_i\pm e_j$ and the half-vectors
$\frac12\sum_i\varepsilon_i e_i$ with $\varepsilon_i=\pm1$
and an even number of minus signs.  Number the simple roots by
\[
\begin{split}
 a_1&=\tfrac12(e_1-e_2-e_3-e_4-e_5-e_6-e_7+e_8),\\
 a_2&=e_1+e_2,\qquad a_3=e_2-e_1,\\
 a_i&=e_{i-1}-e_{i-2}\quad(4\leq i\leq8).
\end{split}
 \tag{EX.6}
\]
Every root has squared length two.  The highest root and its
adjacent root are
\[
 \begin{split}
 \theta&=e_7+e_8
       =2a_1+3a_2+4a_3+6a_4+5a_5+4a_6+3a_7+2a_8,\\
 \delta&=e_6+e_8=\theta-a_8.
 \end{split}
\]
Inspection of the explicit root set shows that $\theta+a_i$
is never a root and that the sole root among $\delta+a_i$ is
$\delta+a_8=\theta$.  In the half-integral case the putative
sum has an entry of absolute value $3/2$; in the integral cases
it has too many nonzero entries or an entry of absolute value two.
The product of reflections in
$e_1-e_2,e_1+e_2,\ldots,e_7-e_8,e_7+e_8$ is $-1$.
The highest-wedge hypotheses therefore hold again.

For the first torus take
$\alpha=e_8-e_1$ and $\beta=\theta$.
Apply the reflections $s_7,s_6,s_5,s_4,s_2$ in that order to
$(\theta,\delta)$.  They all fix $\theta$ and send its other
member successively through
\[
 e_6+e_8,\ e_5+e_8,\ e_4+e_8,\ e_3+e_8,
                    \ e_2+e_8,\ e_8-e_1.
\]
This proves the required pair conjugacy.  The height functional
is $(0,1,2,3,4,5,6,23)$, because it takes value one on every
simple root in (EX.6).  Thus the two heights are $23$ and $29$.

For the second torus put $C=s_1s_2\cdots s_8$ and
$\mu=e_2+e_8$.  The reflection formula
$s_i(v)=v-(v,a_i)a_i$ gives
\[
 C^j\mu=e_{j+2}+e_8\quad(0\leq j\leq5).
\]
The pair $(\mu,C^5\mu)=(e_2+e_8,\theta)$ is obtained from
$(\theta,\delta)$ by $s_7,s_6,s_5,s_4$ in application order.
Here $h=30$, and $5(29-23)=30$.  The obstruction applies with
$d=5$, giving nonopposite containing Borels in type $E_8$.
This proves the claim.  $\square$

\paragraph{Exact supplementary certificate.}
The script \texttt{research/root\_exceptional\_wedge\_check.py}
uses only integers and rational numbers.  It checks weighted Gram
symmetry of the listed Cartan matrices, root counts $48,240$, the raising tests for the
highest wedges, the explicit Weyl words and Coxeter powers, and
the height divisibility.  It also enumerates the highest-pair
orbits, of sizes $96$ and $6720$.  The proofs above need only
the displayed reflection calculations; the enumeration is an
additional finite check, not a numerical substitute for the
highest-line and invariant-pairing argument.

\paragraph{Reference.}
Milne, J. S. (2017). \emph{Algebraic groups: The theory of group
schemes of finite type over a field}. Cambridge University Press.
\url{https://www.jmilne.org/math/Books/iAG2017.pdf}.
Proposition 22.17 and Theorems 22.18--22.19, pp.~467--468;
Proposition 22.41, p.~477; exceptional root systems,
\S24(l), pp.~539--542.  Kostant's exact version and locators
are the same as those already given in CC-APG01--08.
